% Section 7: Classical calibrations: isogenies and the Legendre crossing
\label{sec:classical-body}

A programme that claims a genuinely new triadic structure must be tested against the
places where mathematics already has one.  This section reports two calibrations, and
the first is not flattering to the programme: a complex elliptic isogeny already
carries three exactly-bridged representations, so the bare existence of a bridge is
prior art.  What survives as novel is narrower and is stated precisely.  The second
calibration is constructive: in the Legendre family one cusp really does support three
readings at once, and the residual that obstructs their identification is exact.

\subsection{The isogeny calibration: the bridge is not new}

\begin{theorem}[Three readings of an isogeny]\label{thm:isogeny-three-readings}
Let $\phi\colon E_1\to E_2$ be a complex elliptic isogeny with kernel $\ker\phi$, and let
$M_\alpha$ be the matrix of the induced map on the period lattice attached to a complex
multiplication $\alpha$.  Then
\[
  N(\alpha)=|\det M_\alpha|=\deg\phi=|\ker\phi| ,
\]
so the norm, the lattice determinant, the covering degree, and the kernel cardinality
are four descriptions of one number, computed in three different domains.
\end{theorem}

\begin{proof}
For a complex isogeny the degree equals the size of the kernel and the index of the
image lattice, and the lattice index is the absolute determinant of the matrix of the
map in a lattice basis.  For a CM endomorphism the degree is the algebraic norm.
\end{proof}

\begin{calibration}[The novelty boundary]\label{cal:isogeny-novelty}
Theorem~\ref{thm:isogeny-three-readings} shows that a triadic bridge of the kind
developed in \S\ref{sec:triadic-inference} is already realised classically in the
theory of complex elliptic isogenies.  The programme's claimable novelty is therefore
\emph{not} the existence of three bridged representations.  It is, at most:
\begin{enumerate}
\item observer-relativity: the bridge is stated relative to a declared finite contract
  $Q$ and a hypothesis class;
\item finite inference: the reconstruction is an intersection of fibres subject to a
  budget, with a four-valued outcome;
\item accountable residual: where the classical equalities hold, the residual is
  provably empty, and where they fail, the residual is characterized.
\end{enumerate}
\end{calibration}

\begin{warning}[Two traps in the classical case]\label{warn:isogeny-traps}
Two errors are easy to make and both are avoided in the statement above.
First, coordinate branching of a composite $x\circ\phi$ is \emph{not} ramification of
$\phi$; the two are different structures attached to different maps.
Second, the equalities of Theorem~\ref{thm:isogeny-three-readings} are equalities of
\emph{numbers} obtained from three carriers; they do not identify the carriers, and by
Definition~\ref{def:characteristic} an identification of carriers would require a
residual argument that the theorem does not supply.
\end{warning}

\begin{proposal}[Characteristics as a groupoid]\label{prop:characteristic-groupoid}
A characteristic is better modelled as an object of a representation groupoid --- a
carrier together with declared readings and the isomorphisms between them --- than as a
single normal-form value.  Under this reading, the value of a characteristic is the
isomorphism class, and the residual is the stabilizer datum.
\end{proposal}

\subsection{The marked rank-two carrier and its three cusps}

The geometry that the triadic programme wants is supplied by the level-two modular
curve.

\begin{proposition}[Three cusps are three slope classes]\label{prop:three-cusps}
The level-two modular curve
\[
  Y(2)=\Gamma(2)\backslash\bH\cong\PP^{1}\setminus\{0,1,\infty\}
\]
has exactly three cusps, corresponding to the three points of $\PP^{1}(\F_2)$, that is,
to the three slope classes $0/1$, $1/1$, $1/0$.  Equivalently, they are the three
oriented slope classes of one marked rank-two carrier.
\end{proposition}

\begin{proof}[Status]
The isomorphism is the classical description of $Y(2)$ as the Legendre $\lambda$-line;
the three cusps are $\lambda\in\{0,1,\infty\}$.  The slope-class description is obtained
by reading $(\omega_1,\omega_2)$ modulo $\Gamma(2)$ and reducing the ratio modulo $2$.
The corresponding mod-two relations are
$\delta_T=\delta_S+\delta_C$ and its cyclic variants.
\end{proof}

\begin{remark}[Mod-two erases chirality]\label{rem:mod2-erases}
The reduction above is mod-two, and \S\ref{sec:six-state} showed exactly what
mod-two forgets: the sign of the orientation word.  Consequently the identification of
Proposition~\ref{prop:three-cusps} describes the three readings up to a sign that the
mod-two picture cannot see, and an integer lift is required.  The following
construction supplies it.
\end{remark}

\begin{construction}[Oriented integer lift]\label{con:oriented-lift}
Let $\delta_T=a$, $\delta_S=b$, $\delta_C=-a-b$ in a rank-two lattice with symplectic
form $\langle\cdot,\cdot\rangle$, chosen so that all three cyclic pairings equal $+1$:
\[
  \langle\delta_T,\delta_S\rangle=\langle\delta_S,\delta_C\rangle
  =\langle\delta_C,\delta_T\rangle=+1 .
\]
The order-three chart rotation is
\[
  R=\begin{pmatrix}0&-1\\ 1&-1\end{pmatrix},\qquad R^{3}=I .
\]
\end{construction}

\begin{lemma}[Tri-chart coherence]\label{lem:tri-chart}
The three coordinate charts $(x,y)$, $(y-x,-x)$, $(-y,x-y)$ are permuted by $R$, and
$R^{3}=I$.
\end{lemma}

\begin{proof}
Applying $R$ to a pair yields the next chart, by direct substitution.  The
characteristic polynomial of $R$ is $\lambda^{2}+\lambda+1$, so $R^{3}=I$ by
Cayley--Hamilton.
\end{proof}

\begin{remark}
Construction~\ref{con:oriented-lift} and Lemma~\ref{lem:tri-chart} are the smallest exact
repair of Remark~\ref{rem:mod2-erases}: the sign lost mod two is retained by working in
an oriented integer lattice, at the cost of a nontrivial chart rotation that must be
tracked.
\end{remark}

\subsection{The Legendre crossing}

\begin{proposition}[One cusp, three readings]\label{prop:legendre-three-readings}
Consider the Legendre family
\[
  E_\lambda:\ y^{2}=x(x-1)(x-\lambda),\qquad \lambda\in\PP^{1}\setminus\{0,1,\infty\},
\]
with its nonzero two-torsion satisfying $P_0+P_1+P_\lambda=0$.  At each cusp
$\lambda\in\{0,1,\infty\}$, the same local singularity is simultaneously
\begin{enumerate}
\item a collision of branch factors of the defining equation;
\item a vanishing class $\delta$ with Dehn transvection $U_\delta=L_\delta^{2}$;
\item a monodromy of the Picard--Fuchs system
  \[
    \lambda(1-\lambda)\Pi''+(1-2\lambda)\Pi'-\tfrac14\Pi=0 .
  \]
\end{enumerate}
\end{proposition}

\begin{proof}[Status]
Items (i)--(iii) are classical local descriptions of the Legendre degeneration.  The
content claimed here is only that they apply at one point to one object; the
identification of the three \emph{descriptions} with the three \emph{readings} of
\S\ref{sec:triadic-inference} is not proved and remains the open problem of
\S\ref{sec:boundaries}.
\end{proof}

\begin{nogo}[Bare monodromy is spectrally empty]\label{nogo:unipotent}
The cusp monodromy acts by
\[
  M_\delta(v)=v+\langle v,\delta\rangle\,\delta ,
\]
which is unipotent.  Hence a bare cusp monodromy supplies no rich spectrum, and any
useful spectral invariant must be attached to paths, transports, correspondences, or
gluing data rather than to a single local monodromy operator.
\end{nogo}

\begin{proof}
$M_\delta-\id$ has image contained in the line spanned by $\delta$ and has square zero,
so every eigenvalue of $M_\delta$ equals $1$.
\end{proof}

\begin{proposition}[Central sign and projective closure]\label{prop:central-sign-closure}
With positive unipotent cusp lifts $U_0,U_1,U_\infty$ one has
\[
  U_0U_1U_\infty=-I \quad\text{in the linear lift, while}\quad
  M_0M_1M_\infty=I \quad\text{after typing the sign into the infinity chart.}
\]
Thus the three local operators close projectively, and a central $-1$ survives in any
selected linear lift.
\end{proposition}

\begin{proof}[Status]
The projective closure is classical.  The central sign is the same phenomenon as
Proposition~\ref{prop:order-two-four}: projectivization forgets a scalar, and a lift
retains it.
\end{proof}

\begin{proposition}[Label transport, not universal swap]\label{prop:label-transport}
The general law governing how a vanishing cycle moves when the base chart is changed is
\[
  L_{g\delta}=g\,L_\delta\,g^{-1},
\]
a conjugation of labels.  There is no universal swap law for $\delta$ across charts.
\end{proposition}

\begin{proof}
Conjugation transports the cycle $\delta$ to $g\delta$, and the transvection is
determined by its vanishing cycle, so the conjugated operator is the transvection of the
transported cycle.
\end{proof}

\subsection{The crossing-coherence prism and its decomposable defect}

The three cards are now placed on one object, and the result is the paper's principal
negative calibration.

\begin{construction}[Crossing-coherence prism]\label{con:prism}
Take three typed carriers
\[
  \cK=(\text{construction data }q,p,h),\qquad
  \cX=(\text{charge }q\in\Z^{2}),\qquad
  \cT=(\text{covector }\ell\in(\Z^{2})^{\ast}),
\]
with crossing actions $J^{c}_{\cK}$, $J^{c}_{\cX}=M_c$, $J^{c}_{\cT}=M_c^{-T}$ for each
cusp $c$.  The interpretation maps are $F_{\cK\cX}$, $F_{\cX\cT}=J$, and $F_{\cT\cK}$,
with residuals carried by the last.
\end{construction}

\begin{theorem}[Crossing squares]\label{thm:crossing-squares}
In Construction~\ref{con:prism}:
\begin{enumerate}
\item the $\cK$--$\cX$ square commutes strictly on charge;
\item the $\cX$--$\cT$ square commutes strictly, by symplectic duality
  $JM=M^{-T}J$;
\item the $\cT$--$\cK$ square commutes only up to the exact history residual, recorded
  as $[c]$.
\end{enumerate}
\end{theorem}

\begin{proof}
Item (i) is immediate, since both actions act on the same charge coordinate.  Item (ii)
is the defining property of the symplectic form: for $M$ preserving
$\langle\cdot,\cdot\rangle$ the relation $JM=M^{-T}J$ holds.  Item (iii) fails by
construction, since a period covector does not determine construction history
(Proposition~\ref{prop:period-not-history}); the failure is recorded rather than
removed.
\end{proof}

\begin{proposition}[Two independent residual channels]\label{prop:two-residual-channels}
The prism of Construction~\ref{con:prism} carries two independent residuals:
\begin{enumerate}
\item a central lift sign, of which $-I$ versus $+I$ is the visible instance;
\item a history suffix $[c]$ recording the construction path.
\end{enumerate}
A coarse projective observer that identifies $v\sim-v$ and $M\sim-M$, and that erases the
history, sees neither.
\end{proposition}

\begin{proof}
The first is invisible under projectivization by
Proposition~\ref{prop:order-two-four}; the second is invisible because it is not a
function of the charge or of the covector.  Their independence is exhibited by varying
one while fixing the other.
\end{proof}

\begin{theorem}[No-go: the triadic defect decomposes]\label{thm:defect-decomposes}
The total finite defect of the prism is
\[
  D_Q=\bigl(-I,\,[\infty,1,0]\bigr),
\]
and it decomposes into the two channels of
Proposition~\ref{prop:two-residual-channels}.  Consequently no irreducible triadic
$3$-cocycle is exhibited by this model.
\end{theorem}

\begin{proof}
The central sign is a scalar and the history suffix is a word; the projections onto the
two summands are maps of defect data, and each vanishes on the kernel of the other.  A
defect admitting such a decomposition is by definition not irreducible.
\end{proof}

\begin{calibration}[Status of the negative calibration]\label{cal:defect-negative}
Theorem~\ref{thm:defect-decomposes} is reported as the paper's principal negative
calibration.  It establishes that the finite Legendre prism does \emph{not} produce a
genuinely triadic coherence defect; it does not establish that no such defect exists.
\end{calibration}

\begin{proposal}[What genuine triadicity would require]\label{prop:coherence-defect}
Genuine triadicity requires a coherence defect $\Theta_c$ that is not reconstructible
from any pair of the three crossings.  In the language of
\S\ref{sec:satisfaction-logic} this is the statement that the triple support strictly
exceeds what the three pairwise supports determine.
\end{proposal}

\begin{conjecture}[Triadic singular crossing]\label{conj:triadic-crossing}
There exists a carrier on which the three readings of \S\ref{sec:triadic-inference} are
related by crossing characteristics with a non-decomposable cyclic defect, functorial
under refinement of the observer contract $Q$.
\end{conjecture}

\subsection{What this section settles}

\begin{enumerate}
\item The bare triadic bridge is prior art
  (Theorem~\ref{thm:isogeny-three-readings}, Calibration~\ref{cal:isogeny-novelty}); the
  programme's novelty must be sought in observer-relativity and residual accounting.
\item The classical carrier for three readings is the level-two modular curve with its
  three cusps, and the mod-two picture requires an oriented integer lift to retain
  chirality (Proposition~\ref{prop:three-cusps},
  Construction~\ref{con:oriented-lift}, Lemma~\ref{lem:tri-chart}).
\item Crossing coherence is strict for two pairs and residual-laden for the third, and
  the resulting defect decomposes (Theorem~\ref{thm:crossing-squares},
  Proposition~\ref{prop:two-residual-channels}, Theorem~\ref{thm:defect-decomposes}).
\item Whether a non-decomposable defect exists is open
  (Proposal~\ref{prop:coherence-defect}, Conjecture~\ref{conj:triadic-crossing}).
\end{enumerate}
